\documentclass{article}
\usepackage{amsmath, amssymb, amsthm}
\usepackage{hyperref}
\usepackage{geometry}
\geometry{margin=1in}

\title{Erd\H{o}s Problem \#469}
\date{Last updated: 2025-08-31}
\author{Source: erdosproblems.com}

\begin{document}
\maketitle

\noindent\textbf{Status:} OPEN \quad \textbf{No prize}

\noindent\textbf{Tags:} number theory, divisors

\medskip
\noindent\textbf{Problem Statement:}

Let $A$ be the set of all $n$ such that $n=d_1+\cdots+d_k$ with $d_i$ distinct proper divisors of $n$, but this is not true for any $m\mid n$ with $m<n$. Does\[\sum_{n\in A}\frac{1}{n}\]converge?

\bigskip
\noindent\textbf{Remarks:}

The integers in $A$ are also known as primitive pseudoperfect numbers and are listed as A006036 in the OEIS.

The same question can be asked for those $n$ which do not have distinct sums of sets of divisors, but any proper divisor of $n$ does (which are listed as A119425 in the OEIS).

Benkoski and Erd\H{o}s \cite{BeEr74} ask about these two sets, and also about the set of $n$ that have a divisor expressible as a distinct sum of other divisors of $n$, but where no proper divisor of $n$ has this property.

\bigskip
\noindent\textbf{References:}

[BeEr74] Benkoski, S. J. and Erd\H{o}s, P., On weird and pseudoperfect numbers. Math. Comp. (1974), 617-623.

\bigskip
\noindent\small{Source: \url{https://www.erdosproblems.com/469}}
\end{document}
